\begin{lemma}[Determinism of the Minimal Iterator Relation]
\label{lem:minimal-iterator-relation-deterministic}
Let $(P,S,1)$ be a Peano system, let $(W,c,g)$ be iterator data for it, and let
$R^*$ be the minimal iterator relation. For each $n\in P$, there is at most one
$w\in W$ such that $(n,w)\in R^*$.
\hyperref[prf:minimal-iterator-relation-deterministic]{Go to proof.}
\end{lemma}
\end{lemmabox}

\begin{formalizationrecord}{Lean 4 Verification Record}
\FormalizationSource{Feferman, \emph{The Number Systems}, Section 3.4; Mendelson, \emph{Number Systems and the Foundations of Analysis}, Section 2.2}
\Formalizes{lem:minimal-iterator-relation-deterministic}
\textbf{Lean module:} \texttt{LRA.VolumeII.PeanoSystems.Recursion.Iterator}.\par
\LeanDeclaration{MinimalIteratorRelationDeterministic}
\textbf{Status:} pending against \texttt{lra-lean}.\par

\begin{leancode}
theorem MinimalIteratorRelationDeterministic
    (ps : PeanoSystem)
    (data : IteratorData ps) :
    forall element : ps.carrier,
      forall first_value second_value : data.Target,
        MinimalIteratorRelation ps data element first_value ->
        MinimalIteratorRelation ps data element second_value ->
        first_value = second_value := by
  sorry
\end{leancode}
\end{formalizationrecord}

\begin{remark*}[Standard quantified statement]
\[
\forall n\in P\;\forall w_0,w_1\in W\;
\bigl((n,w_0)\in R^*\land (n,w_1)\in R^*\Longrightarrow w_0=w_1\bigr).
\]
\end{remark*}

\begin{remark*}[Predicate reading]
\[
\operatorname{FunctionalRelation}(R^*,P,W).
\]
\end{remark*}

\begin{remark*}[Negated quantified statement]
\[
\exists n\in P\;\exists w_0,w_1\in W\;
\bigl((n,w_0)\in R^*\land (n,w_1)\in R^*\land w_0\neq w_1\bigr).
\]
\end{remark*}

\begin{remark*}[Failure modes]
\begin{description}
  \item[Exposition.]
  Determinism fails exactly when a single stage has two distinct values in the
  minimal iterator relation.
  \item[Two values.]
  \[
  \neg\operatorname{FunctionalRelation}(R^*,P,W).
  \]
  \[
  \exists n\in P\;\exists w_0,w_1\in W\;
  \bigl((n,w_0)\in R^*\land (n,w_1)\in R^*\land w_0\neq w_1\bigr).
  \]
\end{description}
\end{remark*}

\begin{remark*}[Interpretation]
Determinism says the minimal iterator relation is functional at every stage.
This repackages forced-value uniqueness in the at-most-one form needed to
extract a function from the relation.
\end{remark*}

\begin{dependencies}
\begin{itemize}
  \item \hyperref[lem:forced-values-are-unique]{Forced Values Are Unique}
\end{itemize}
\end{dependencies}

